\section{Conclusion}
\label{sec:conclusion}

We studied how equivalent creative-writing instructions change model behavior when supplied upfront or delivered progressively across a conversation (\cref{sec:methodology}). Across 160 tasks, six models, and 1{,}920 final generations, FULL prompting produced substantially higher constraint adherence and stronger automated craft, genre effectiveness, and especially structure/coherence (\cref{sec:rq1,sec:rq2}). Importantly, the structure/coherence advantage persisted among pairs with equal observed constraint adherence (\cref{sec:rq3}), suggesting that multi-turn degradation is not explained solely by forgetting explicit requirements. At the same time, Sharded outputs received slightly higher automated originality scores, indicating that the effect is not a uniform decline across all writing qualities. A 30-pair blinded human study further showed that LLM judgments aligned considerably better with humans on constraint following than on holistic creative preference (\cref{sec:disc_human}). Together, these results suggest that progressive specification affects both how models maintain narrative state and how reliably those changes can be evaluated.

Future work can extend the human evaluation, analyze turn-level trajectories and constraint position, and test whether the observed trade-offs persist in larger models and longer conversations.
